%%%%%
%
\section{Conversion d'une angle donn\'e entre 0 et 360 degr\'es sur la gamme [-180;+180]}

\itemstyle{Op\'erateur~:} {\sf dd2pi2ddpi}
\index{dd2pi2ddpi}

%%%%%OBJET%%%%
\itemstyle{Objet~:}
  Convertit un angle compt\'e entre 0 et 360 degr\'es en un angle variant de -180 \`a +180 degr\'es.

%%%%%ARGUMENTS%%%%
\bigskip
%%%%%ENTREES%%%%
\begin{tabarg}
\multicolumn{3}{l}{\titretab {Entr\'ees}}\\
\hline
$dda2$ & Vecteur, \texttt {float} & Angle sur [0;360]  (en degr\'es) \\
\end{tabarg}

\medskip
%%%%%SORTIES%%%%
\begin{tabarg}
\multicolumn{3}{l}{\titretab {Sorties}}\\
\hline
$dda1$ & Vecteur, \texttt {float} & Angle sur [-180;180] (en degr\'es) \\
\end{tabarg}



%%%%%FORMULE%%%%
\itemstyle{Formule}

%%dd2pi2ddpi 
%\usepackage{amsmath}
\begin{equation*}
\begin{aligned}
dda1 = &
\begin{cases}
dda2-360 & \text{si $dda2 > 180$}\\
dda2 & \text{si $dda2 <= 180$}\\
\end{cases} \\
\end{aligned}
\end{equation*}



%%%%%COMMENTAIRE%%%%
\itemstyle{Commentaire~:}
Il est utile pour certains op\'erateurs de changement de rep\`eres, et en particuliers pour les calculs de vents vrais.
Dans ce cas, dda2 peut \^etre par exemple, le vent apparent mesur\'e sur un navire.
   
%%%%%AUTEUR%%%%
\itemstyle{Auteur~:}
CNRM/GMEI/TRAMM

%%%%%REFERENCE%%%%
\itemstyle{R\'ef\'erences~:}
